\chapter{条件制备、计数与统计反演}\label{chap:conditional-electron-inference}

同一个联合作用既可以给出无条件电子统计，也可以在读取某个结果后制备光场。条件操作、有限计数核与反演因此要使用相容的归一化和同一前向模型。

\section{条件测量与统计反演}

电子与光相互作用后通常存在相关。若测量电子并只保留某一个结果，与这些电子对应的光场统计也会改变。这与条件概率的思想相同，但光场还具有相干，因此需要更新整个密度算符。本节由联合散射和电子读出求条件光场态，并讨论怎样从多次电子计数推断入射光子统计。

\subsection{条件光场仪器}

现在把光场视为待保留的系统，把电子视为探针。先由联合散射算符计算电子与光的关联，再对电子施加前节定义的测量效应。对电子求偏迹后，得到光场的未归一化条件态；它的迹是电子结果的概率，除以该概率才是相应的光场态。这就从已有的联合动力学构造出作用在光场上的量子仪器。

设电子入射态在能量格基底中为
\begin{equation*}
\hat\rho_e
=\sum_{j,j'}r_{jj'}|j\rangle\langle j'|,
\qquad
\sum_j r_{jj}=1.
\end{equation*}
$r_{jj'}$ 无量纲；$r_{jj}$ 是电子初始能量格占据概率，$r_{jj'}$ 的非对角元描述电子能量相干。设光场初态为 $\hat\rho_{\rm ph}$，联合散射算符为无量纲幺正算符 $\hat S_Q$。

若电子理想测得末态格点 $\ell$，对每一个初始电子格点 $j$ 定义作用在光场希尔伯特空间上的\emph{条件偏矩阵元}
\begin{equation*}
\mathscr M_{\ell j}
\equiv {}_e\!\langle\ell|\hat S_Q|j\rangle_e.
\end{equation*}
$\mathscr M_{\ell j}$ 是无量纲光场算符；电子指标 $\ell,j$ 只标记取了哪个电子矩阵元，并不是光场算符的矩阵指标。理想结果 $\ell$ 对应的未归一化条件光场态为
\begin{equation}
\boxed{
\widetilde\rho_{\rm ph|\ell}
=\sum_{j,j'}r_{jj'}
\mathscr M_{\ell j}\hat\rho_{\rm ph}
\mathscr M_{\ell j'}^\dagger,
\qquad
P_\ell
=\operatorname{Tr}_{\rm ph}\widetilde\rho_{\rm ph|\ell}.}
\label{eq:general_electron_meter_conditional_photon_map}
\end{equation}
$\widetilde\rho_{\rm ph|\ell}$ 无量纲但迹一般小于一；其迹 $P_\ell$ 是无量纲测量概率。若 $P_\ell>0$，归一化条件态为 $\rho_{\rm ph|\ell}=\widetilde\rho_{\rm ph|\ell}/P_\ell$。

电子测量在光场上诱导的效应算符为
\begin{equation*}
\hat\Pi_\ell^{({\rm ph})}
\equiv\sum_{j,j'}r_{jj'}
\mathscr M_{\ell j'}^\dagger\mathscr M_{\ell j},
\qquad
P_\ell=\operatorname{Tr}_{\rm ph}
[\hat\Pi_\ell^{({\rm ph})}\rho_{\rm ph}].
\end{equation*}
联合幺正性与电子末态完备性保证
\begin{equation*}
\sum_\ell \hat\Pi_\ell^{({\rm ph})}=\mathbb I_{\rm ph}.
\end{equation*}
因此 $\{\hat\Pi_\ell^{({\rm ph})}\}$ 是作用在光场上的 POVM。这里的 $\hat\Pi_\ell^{({\rm ph})}$ 只是测量效应算符的统一记号，并不要求它满足投影条件 $\hat\Pi^2=\hat\Pi$；只有理想投影测量才具有这一额外性质。

QPINEM 单模交换还满足守恒律\eqnrefs{eq:qpinem_conservation}。用已经冻结的总交换量
\[
\mathcal N_{\rm ex}=\Lambda_e+N_{\rm ph}
\]
表示电子能量格阶数与光子数之和，则电子从 $j$ 变到 $\ell$ 时，光子数必须由 $n$ 变到 $n+j-\ell$。因此 $r_{jj'}$ 的第 $j-j'$ 阶电子相干会把光场相隔 $j-j'$ 个 Fock 格点的相干带入同一个电子计数概率。单能电子入射时 $r_{jj'}=\delta_{j0}\delta_{j'0}$，这些交叉项消失，光场 POVM 才退化为光子数对角形式。

真实能谱仪还具有有限效率和有限能量分辨率。前面定义的探测器响应密度
\[
R_{\rm det}(E_d|\mathcal E_\ell)
=\mathcal Q_\ell^{\rm det} r(E_d|\mathcal E_\ell)
\]
单位为 $\mathrm{J^{-1}}$，并满足 $\int\dd E_d\,R_{\rm det}=\mathcal Q_\ell^{\rm det}$。因此记录连续能量结果 $E_d$ 后的\emph{算符值概率密度}为
\begin{equation}
\boxed{
\widetilde\rho_{\rm ph|E_d}
=\sum_\ell R_{\rm det}(E_d|\mathcal E_\ell)
\sum_{j,j'}r_{jj'}
\mathscr M_{\ell j}\hat\rho_{\rm ph}
\mathscr M_{\ell j'}^\dagger.}
\label{eq:conditional_photon_detector_coarse_graining}
\end{equation}
因为 $R_{\rm det}$ 带 $\mathrm{J^{-1}}$，所以
\[
[\widetilde\rho_{\rm ph|E_d}]=\mathrm{J^{-1}},
\]
并且 $\operatorname{Tr}\widetilde\rho_{\rm ph|E_d}\,\dd E_d$ 是在区间 $[E_d,E_d+\dd E_d]$ 内得到点击的概率。若实验只接受能窗 $\mathcal W$，定义无量纲权重
\[
w_\ell(\mathcal W)
\equiv\int_{\mathcal W}\dd E_d\,
R_{\rm det}(E_d|\mathcal E_\ell),
\qquad 0\le w_\ell\le\mathcal Q_\ell^{\rm det},
\]
即可得到窗口后选择的条件量子仪器。

\subsection{单模计数核与条件仪器}\label{subsec:single-mode-number-instrument-dp19}

在单模 QPINEM 中，前面已经求出每个光子数输入对应的联合散射振幅。若只统计电子落在哪条边带，对不同的未测光子末态求和即可得到概率。若希望在电子后选择后继续使用光场，则要保留这些振幅的相位，并用 Kraus 算符表示条件更新。下面从同一组散射矩阵元写出这两种结果。

\begin{exbox}[QPINEM 单模数统计]
联合散射给出每个初始光子数 $n$ 到电子边带 $\ell$ 的振幅。若只读取电子边带而不保留末态光场相位，所需信息只是这些振幅的模平方，因此形成从光子数概率到电子计数概率的随机核；若进一步按某个电子结果后选择，光场振幅本身仍被保留，于是得到作用在光场上的条件算符。二者都由同一联合散射矩阵元导出。

取入射电子为单一能量格点 $|0\rangle_e$，光场密度算符写成
\begin{equation*}
\hat\rho_{\rm ph}
=\sum_{n,n'\ge0}\rho_{nn'}|n\rangle\langle n'|,
\qquad
p_n\equiv\rho_{nn}.
\end{equation*}
$p_n$ 是无量纲光子数概率，满足 $p_n\ge0$ 与 $\sum_n p_n=1$。

电子测得边带 $\ell$ 时，守恒量把初始光子数 $n$ 与末态光子数 $n-\ell$ 绑定。对末态光场求迹后，不同 Fock 态正交，因此电子边带概率只依赖光子数分布。定义无量纲随机核
\begin{equation}
\boxed{
P_\ell
=\sum_{n=0}^{\infty}\mathsf K_{\ell n}(|g_{\mathrm Q}|)p_n,
\qquad
\mathsf K_{\ell n}(|g_{\mathrm Q}|)
\equiv|\mathcal A_{\ell|n}(g_{\mathrm Q})|^2.}
\label{eq:electron_spectrum_linear_kernel}
\end{equation}
$\mathsf K_{\ell n}$ 是条件概率 $P(\ell|n)$，因此
\begin{equation*}
\mathsf K_{\ell n}\ge0,
\qquad
\sum_\ell\mathsf K_{\ell n}=1
\end{equation*}
对每个固定 $n$ 成立。这里 $g_{\mathrm Q}$ 是 QPINEM 单模无量纲交换振幅。由于 $\mathcal A_{\ell|n}(g_{\mathrm Q})$ 来自该单模有效模型的完整位移算符矩阵元，$\mathsf K_{\ell n}$ 包含 $|g_{\mathrm Q}|$ 的全部阶次；弱耦合响应是这一模型内全阶核在 $|g_{\mathrm Q}|\ll1$ 时的 Taylor 展开。

由位移算符矩阵元~\eqnrefs{eq:A_l_n_displacement}，增益边带 $\ell\ge0$ 的模型内全阶核为
\begin{equation}
\boxed{
\mathsf K_{\ell n}
=\ee^{-|g_{\mathrm Q}|^2}\frac{(n-\ell)!}{n!}|g_{\mathrm Q}|^{2\ell}
\left[L_{n-\ell}^{(\ell)}(|g_{\mathrm Q}|^2)\right]^2,
\qquad n\ge\ell,}
\label{eq:kernel_gain}
\end{equation}
而电子损失 $r_{\rm loss}=-\ell\ge0$ 时
\begin{equation}
\boxed{
\mathsf K_{-r_{\rm loss},n}
=\ee^{-|g_{\mathrm Q}|^2}\frac{n!}{(n+r_{\rm loss})!}|g_{\mathrm Q}|^{2r_{\rm loss}}
\left[L_n^{(r_{\rm loss})}(|g_{\mathrm Q}|^2)\right]^2.}
\label{eq:kernel_loss}
\end{equation}

与随机核不同，理想电子结果 $\ell$ 对应的光场条件算符定义为
\begin{equation*}
\mathscr M_\ell
\equiv{}_e\!\langle\ell|\hat S_Q|0\rangle_e
=\sum_{n\ge\max(0,\ell)}
\mathcal A_{\ell|n}|n-\ell\rangle\langle n|.
\end{equation*}
$\mathscr M_\ell$ 无量纲，作用于光场 Hilbert 空间。相应 POVM 效应和归一化条件态为
\begin{equation}
\boxed{
\hat\Pi_\ell^{({\rm ph})}
=\mathscr M_\ell^\dagger\mathscr M_\ell
=\sum_n\mathsf K_{\ell n}|n\rangle\langle n|,
\qquad
\rho_{\rm ph|\ell}
=\frac{\mathscr M_\ell\rho_{\rm ph}\mathscr M_\ell^\dagger}{P_\ell}.}
\label{eq:single_electron_photon_instrument}
\end{equation}
联合幺正性给
\[
\sum_\ell \hat\Pi_\ell^{({\rm ph})}=\mathbb I_{\rm ph},
\]
所以 $\sum_\ell P_\ell=1$。

弱交换极限把上述模型内全阶核与光子下降阶乘矩连接起来。定义
\begin{equation*}
F_{r_{\rm ph}}
\equiv\left\langle
(a^\dagger)^{r_{\rm ph}}a^{r_{\rm ph}}
\right\rangle,
\qquad r_{\rm ph}\in\mathbb N_+,
\qquad
[F_{r_{\rm ph}}]=1.
\end{equation*}
$F_{r_{\rm ph}}$ 是无量纲第 $r_{\rm ph}$ 阶下降阶乘矩。若平均光子数 $\bar n=\langle N\rangle>0$，归一化零延迟关联定义为
\begin{equation*}
g_{\rm ph}^{(r_{\rm ph})}(0)
\equiv\frac{F_{r_{\rm ph}}}{\bar n^{r_{\rm ph}}}.
\end{equation*}
由弱交换展开，电子增益第 $r_{\rm ph}$ 阶边带满足\eqnrefs{eq:qpinem-weak-factorial-moment-dp19}，从而
\begin{equation}
\boxed{
g_{\rm ph}^{(r_{\rm ph})}(0)
=\lim_{|g_{\mathrm Q}|\to0}
\frac{(r_{\rm ph}!)^2P_{+r_{\rm ph}}}{P_{+1}^{r_{\rm ph}}}.}
\label{eq:electron_factorial_correlation_readout}
\end{equation}
有限 $|g_{\mathrm Q}|$ 时应直接使用 $\mathsf K_{\ell n}$，因为高阶 $|g_{\mathrm Q}|$ 修正会同时混合多个光子数矩。

\medskip

若只读出电子能谱，光场信息通过联合散射后电子结果的条件概率进入。该关系可写成从光子统计到电子计数的随机核。 从联合散射和 Born 规则推出电子谱线性核~\eqnrefs{eq:electron_spectrum_linear_kernel}，并重写为单电子对光场的 instrument~\eqnrefs{eq:single_electron_photon_instrument}。 这里重复电子事件独立；输入光场态在单次事件内固定；核的列归一化由概率守恒保证。对初态 $|0\rangle_e\langle0|\otimes\rho_{\rm ph}$，先展开光场密度算符：
\[
\rho_{\rm ph}=\sum_{n,n'}\rho_{nn'}|n\rangle\langle n'|.
\]
联合散射作用在右矢上：
\[
\hat S_Q|0,n\rangle
=\sum_\ell\mathcal A_{\ell|n}|\ell,n-\ell\rangle.
\]
左矢对应共轭展开。投影电子结果 $\ell$ 后，概率为
\begin{align*}
P_\ell
&=\sum_{n,n'}\rho_{nn'}
\mathcal A_{\ell|n}\mathcal A_{\ell|n'}^*
\langle n'-\ell|n-\ell\rangle\\
&=\sum_{n,n'}\rho_{nn'}
\mathcal A_{\ell|n}\mathcal A_{\ell|n'}^*\delta_{nn'}\\
&=\sum_n p_n|\mathcal A_{\ell|n}|^2.
\end{align*}
因此 $\mathsf K_{\ell n}=|\mathcal A_{\ell|n}|^2$。对固定初始 $n$，联合幺正性要求
\begin{align*}
1
&=\langle0,n|\hat S_Q^\dagger\hat S_Q|0,n\rangle\\
&=\sum_\ell|\mathcal A_{\ell|n}|^2
=\sum_\ell\mathsf K_{\ell n},
\end{align*}
所以每一列 $\mathsf K_{\ell n}$ 都是归一化条件概率分布。保留末态光场而只对电子取偏矩阵元：
\[
\mathscr M_\ell
={}_e\!\langle\ell|\hat S_Q|0\rangle_e.
\]
作用在 $|n\rangle$ 上得到
\[
\mathscr M_\ell|n\rangle
=\mathcal A_{\ell|n}|n-\ell\rangle,
\]
所以
\[
\mathscr M_\ell
=\sum_n\mathcal A_{\ell|n}|n-\ell\rangle\langle n|.
\]
于是
\begin{align*}
\mathscr M_\ell^\dagger\mathscr M_\ell
&=\sum_{n,n'}
\mathcal A_{\ell|n}^*\mathcal A_{\ell|n'}
|n\rangle\langle n-\ell|n'-\ell\rangle\langle n'|\\
&=\sum_n|\mathcal A_{\ell|n}|^2|n\rangle\langle n|\\
&=\sum_n\mathsf K_{\ell n}|n\rangle\langle n|.
\end{align*}
再用电子末态完备性：
\begin{align*}
\sum_\ell\mathscr M_\ell^\dagger\mathscr M_\ell
&={}_e\!\langle0|\hat S_Q^\dagger
\left(\sum_\ell|\ell\rangle\langle\ell|\right)
\hat S_Q|0\rangle_e\\
&={}_e\!\langle0|\hat S_Q^\dagger\hat S_Q|0\rangle_e\\
&=\mathbb I_{\rm ph}.
\end{align*}

最后，把精确 Laguerre 振幅\eqnrefs{eq:A_l_n_displacement} 取模平方。对于 $\ell\ge0$，相位因子模长为一，剩余阶乘与 Laguerre 多项式直接给\eqnrefs{eq:kernel_gain}；对于 $\ell=-r_{\rm loss}\le0$，对应另一支矩阵元给\eqnrefs{eq:kernel_loss}。
\end{exbox}

\subsection{有限计数核的反演与条件操作}\label{subsec:statistics-inversion-conditioning-dp19}

设只需重建有限个光子数概率，电子探测器也只读取有限个能量区间，那么前面的计数关系成为一个有限维矩阵方程。要从电子计数求出光子概率，矩阵的不同列必须提供独立信息；若两列几乎相同，少量计数误差就会造成很大的反演误差。我们先用秩和奇异值检查这一点，再计算选定电子结果后的光场更新。

\begin{exbox}[光子统计反演与电子后选择条件操作]
\eqnrefs{eq:electron_spectrum_linear_kernel} 把有限单模光子数统计映射为电子边带分布。Fock 态、相干态和热态共享同一个条件响应核 $\mathsf K_{\ell n}$，它们的差异由输入光子数分布 $p_n$ 进入。典型状态的下降阶乘矩为
\begin{equation*}
\begin{array}{c|c|c}
\text{光态}&F_{r_{\rm ph}}&g_{\rm ph}^{(r_{\rm ph})}(0)\\ \hline
|N\rangle&N(N-1)\cdots(N-r_{\rm ph}+1)&\dfrac{N!}{(N-r_{\rm ph})!N^{r_{\rm ph}}}\\[4pt]
|\alpha_c\rangle&\bar n^{r_{\rm ph}}&1\\[2pt]
\rho_{\rm th}&r_{\rm ph}!\,\bar n^{r_{\rm ph}}&r_{\rm ph}!
\end{array}
\end{equation*}
这里用 $\alpha_c$ 表示相干态振幅，避免与精细结构常数 $\alpha$ 混用；$\bar n=|\alpha_c|^2$ 或热平均光子数均无量纲。

热态还存在一个受控的经典标度极限。令
\begin{equation*}
|g_{\mathrm Q}|\to0,
\qquad
\bar n\to\infty,
\qquad
\bar N_{\rm ex}^{(\rm th)}\equiv|g_{\mathrm Q}|^2\bar n
\end{equation*}
保持有限，则 $\bar N_{\rm ex}^{(\rm th)}$ 是无量纲平均交换强度。此时电子边带趋于
\begin{equation}
\boxed{
P_\ell^{({\rm th,cl})}
=\ee^{-2\bar N_{\rm ex}^{(\rm th)}}I_{|\ell|}(2\bar N_{\rm ex}^{(\rm th)}),
\qquad
\operatorname{Var}(\ell)=2\bar N_{\rm ex}^{(\rm th)}.}
\label{eq:thermal_skellam}
\end{equation}
该分布是对称 Skellam 型离散分布；它描述弱单光子耦合与大热占据联合缩放下的极限，只描述上述联合标度；有限 $g_{\mathrm Q}$ 时应使用单模有效模型的全阶 Laguerre 核。

实验反演只能使用有限电子边带集合 $\mathcal L_{\rm det}$，并通常把待估光子数截断到 $0\le n\le N_{\max}$。把观测边带概率组成向量 $\mathbf P$，把未知光子数分布组成 $\mathbf p$，定义截断后的无量纲设计矩阵
\begin{equation}
\boxed{
\mathsf K_{\rm inv}
\equiv[\mathsf K_{\ell n}]_{\ell\in\mathcal L_{\rm det},\,0\le n\le N_{\max}},
\qquad
\mathbf P=\mathsf K_{\rm inv}\mathbf p.}
\label{eq:single_mode_kernel_inverse_problem_dp68}
\end{equation}
这里 $\mathsf K_{\rm inv}$ 是从单模有效模型的全阶前向核中抽取出的有限数据矩阵。其奇异值分解
\begin{equation}
\boxed{
\mathsf K_{\rm inv}
=\mathsf U_K\operatorname{diag}(s_1^{(K)},s_2^{(K)},\ldots)\mathsf V_K^\dagger}
\label{eq:single_mode_kernel_svd}
\end{equation}
直接给出反演可辨识性。由于概率还满足 $\sum_n p_n=1$，真正的零空间与条件数应限制在归一化约束的切空间上。若 $s_a^{(K)}=0$，相应光子数方向不改变当前电子数据；若 $s_a^{(K)}\ll1$，伪逆沿该方向把数据噪声放大约 $1/s_a^{(K)}$。改变 $|g_{\mathrm Q}|$、电子参考相位或预整形设置相当于增加新的前向矩阵块，可以改变联合奇异值谱。

同一个联合散射还可以用来制备条件光态。设入射电子为能量梳
\begin{equation*}
|\psi_e\rangle
=\sum_j c_j|j\rangle,
\qquad
\sum_j|c_j|^2=1,
\end{equation*}
其中 $c_j$ 是无量纲复振幅。后选择电子末态 $m$ 后，定义无量纲条件光场算符
\begin{equation}
\boxed{
\mathscr C_m[\psi_e]
\equiv{}_e\!\langle m|\hat S_Q|\psi_e\rangle_e
=\sum_j c_j\,\mathscr M_{mj}.}
\label{eq:comb_conditioned_kraus-dp19}
\end{equation}
若平移协变性使 $\mathscr M_{mj}$ 只依赖净电子位移 $m-j$，则可进一步写成 $\sum_jc_j\mathscr M_{m-j}$。因此电子梳的相对相位会相干叠加不同光子增减路径，条件操作取决于振幅而不只取决于电子能量概率。

有限宽电子相位参考态 $|\theta;M\rangle$ 由 $2M+1$ 个能量格点构成。其对双边平移算符的边界误差来源于有限格宽；随着 $M$ 增大，归一化边界振幅按 $M^{-1/2}$ 衰减，而相位主瓣宽度按 $M^{-1}$ 缩小。这两个尺度分别控制“近似平移本征态”误差和可用相位分辨率。

\medskip

热光场下吸收与发射事件在经典大占据数极限可近似为两个独立 Poisson 过程，其净边带阶数是二者之差。 由两个 Poisson 生成函数的乘积推出 Skellam 分布~\eqnrefs{eq:thermal_skellam}，并给出均值与方差；这里经典标度极限使单次跃迁稀疏而总事件数有限；正负跃迁独立。热态的 Glauber--Sudarshan 表示为
\[
\rho_{\rm th}
=\int\frac{\dd^2 \alpha_c}{\pi\bar n}
\exp\!\left(-\frac{|\alpha_c|^2}{\bar n}\right)
|\alpha_c\rangle\langle\alpha_c|.
\]
在联合标度 $|g_{\mathrm Q}|\to0$、$\bar n\to\infty$ 且 $\bar N_{\rm ex}^{(\rm th)}=|g_{\mathrm Q}|^2\bar n$ 固定时，给定相干振幅 $\alpha_c$ 的电子边带趋于经典 Bessel 核
\[
P_\ell(\alpha_c)=J_\ell^2(2|g_{\mathrm Q}\alpha_c|).
\]
写 $\alpha_c=re^{\ii\varphi}$。核只依赖 $r$，所以角变量积分给 $2\pi$：
\begin{align*}
P_\ell
&=\int_0^\infty\frac{2r\dd r}{\bar n}
\exp\!\left(-\frac{r^2}{\bar n}\right)
J_\ell^2(2|g_{\mathrm Q}|r).
\end{align*}
令
\[
x=\frac{r^2}{\bar n},
\qquad
r\dd r=\frac{\bar n}{2}\dd x,
\]
得到
\[
P_\ell
=\int_0^\infty\dd x\,\ee^{-x}
J_\ell^2(2\sqrt{\bar N_{\rm ex}^{(\rm th)}x}).
\]
与其直接调用积分表，先构造边带特征函数
\[
\chi_{\rm sb}(\vartheta)
\equiv\sum_{\ell=-\infty}^{\infty}
P_\ell\ee^{\ii\ell\vartheta}.
\]
利用 Bessel 加法恒等式
\[
\sum_{\ell=-\infty}^{\infty}
J_\ell^2(z)\ee^{\ii\ell\vartheta}
=J_0\!\left(2z\sin\frac{\vartheta}{2}\right),
\]
有
\[
\chi_{\rm sb}(\vartheta)
=\int_0^\infty\dd x\,\ee^{-x}
J_0\!\left(4\sqrt{\bar N_{\rm ex}^{(\rm th)}x}\sin\frac{\vartheta}{2}\right).
\]
把 $J_0$ 展开成幂级数
\[
J_0(2\sqrt{ax})
=\sum_{m=0}^\infty
\frac{(-a)^m x^m}{(m!)^2},
\]
其中此处
\[
a=4\bar N_{\rm ex}^{(\rm th)}\sin^2\frac{\vartheta}{2}.
\]
逐项积分并使用
\[
\int_0^\infty\dd x\,\ee^{-x}x^m=m!,
\]
得到
\begin{align*}
\chi_{\rm sb}(\vartheta)
&=\sum_{m=0}^\infty\frac{(-a)^m}{m!}\\
&=\ee^{-a}\\
&=\exp\!\left[-4\bar N_{\rm ex}^{(\rm th)}\sin^2\frac{\vartheta}{2}\right]\\
&=\exp[2\bar N_{\rm ex}^{(\rm th)}(\cos\vartheta-1)].
\end{align*}
修正 Bessel 函数生成式为
\[
\exp[x(t+t^{-1})]
=\sum_{\ell=-\infty}^{\infty}I_{|\ell|}(2x)t^\ell.
\]
取 $t=\ee^{\ii\vartheta}$，
\[
\exp[2x\cos\vartheta]
=\sum_\ell I_{|\ell|}(2x)\ee^{\ii\ell\vartheta}.
\]
因此
\[
\chi_{\rm sb}(\vartheta)
=\sum_\ell
\ee^{-2\bar N_{\rm ex}^{(\rm th)}}I_{|\ell|}(2\bar N_{\rm ex}^{(\rm th)})
\ee^{\ii\ell\vartheta},
\]
逐项比较 Fourier 系数即得到\eqnrefs{eq:thermal_skellam}。由于
\[
\langle\ell\rangle
=\frac1{\ii}\chi_{\rm sb}'(0)=0,
\]
并且
\[
\langle\ell^2\rangle
=-\chi_{\rm sb}''(0)=2\bar N_{\rm ex}^{(\rm th)},
\]
所以方差为 $2\bar N_{\rm ex}^{(\rm th)}$。

\medskip

\eqnrefs{eq:single_mode_kernel_inverse_problem_dp68} 给出有限维线性前向问题
\[
\mathbf P=\mathsf K_{\rm inv}\mathbf p.
\]
对实测矩阵作奇异值分解
\[
\mathsf K_{\rm inv}=\mathsf U_K\mathsf S_K\mathsf V_K^\dagger,
\qquad
\mathsf S_K=\operatorname{diag}(s_1^{(K)},s_2^{(K)},\ldots),
\]
其中 $\mathsf U_K,\mathsf V_K$ 酉。把未知分布的微小变化展开在右奇异向量 $\mathbf v_a$ 上，
\[
\delta\mathbf p=\sum_a c_a\mathbf v_a,
\qquad
\mathsf V_K^\dagger\mathbf v_a=\mathbf e_a,
\]
则数据变化为
\begin{align*}
\delta\mathbf P
&=\mathsf K_{\rm inv}\delta\mathbf p\\
&=\mathsf U_K\mathsf S_K\mathsf V_K^\dagger
\sum_a c_a\mathbf v_a\\
&=\sum_a c_a s_a^{(K)}\mathbf u_a.
\end{align*}
因此 $s_a^{(K)}=0$ 时，沿 $\mathbf v_a$ 的参数变化严格落在前向映射零空间中；若 $s_a^{(K)}>0$，形式伪逆对相应数据分量给
\[
\delta\mathbf p_a
=\frac{\mathbf v_a\mathbf u_a^\dagger\delta\mathbf P}{s_a^{(K)}}.
\]
所以该方向上的数据噪声被放大 $1/s_a^{(K)}$。概率归一化要求 $\mathbf 1^{\mathsf T}\delta\mathbf p=0$，实际可辨识性应在这一归一化切空间内检查，即正文对零空间和条件数的限定。

\medskip

\eqnrefs{eq:general_electron_meter_conditional_photon_map} 已说明任意入射电子密度矩阵如何通过联合散射的电子偏矩阵元条件化光场。若入射电子进一步取纯的能量梳态，条件映射在振幅层可以直接重写为这些偏矩阵元的相干叠加；这一步给出\eqnrefs{eq:comb_conditioned_kraus-dp19}，并在平移协变时化成卷积。

对电子能量梳，\begin{align*}
\mathscr C_m[\psi_e]
&={}_e\!\langle m|\hat S_Q|\psi_e\rangle_e\\
&={}_e\!\langle m|\hat S_Q
\sum_jc_j|j\rangle_e\\
&=\sum_jc_j\,{}_e\!\langle m|\hat S_Q|j\rangle_e\\
&=\sum_jc_j\mathscr M_{mj}.
\end{align*}
若电子格平移协变，存在单位平移算符 $T_e$ 使 $|j\rangle=T_e^j|0\rangle$，并且联合散射只依赖净电子格位移。于是
\[
\mathscr M_{mj}=\mathscr M_{m-j},
\]
从而得到卷积形式
\[
\mathscr C_m[\psi_e]
=\sum_jc_j\mathscr M_{m-j}.
\]
每一项保留 $c_j$ 的复相位，所以能量梳相位会在条件光场振幅层发生相干干涉，而不是在概率层作非相干求和。
\end{exbox}


\subsection{聚束最优传播、能散与退相干}\label{subsec:bunching-decoherence-closure-dp47}

\eqnrefs{eq:electron_bunching_harmonic_dp46} 已把散射后能量格密度矩阵映射为传播后的密度谐波。实验上真正需要的是在给定漂移距离、输入能散和环境噪声下的可见谐波。令 $t=L/v_0$，并把理想聚束因子记为 $B_{\Delta\ell}^{(0)}(L)$。若输入电子具有小的速度涨落 $\delta v$，则同一实验距离对应传播时间涨落
\begin{equation*}
 \delta t\simeq-\frac{L}{v_0^2}\delta v.
\end{equation*}
对窄带相对论电子，$\delta v=(\partial v/\partial E)_{E_0}\delta E$。定义
\begin{equation}
\boxed{
 \sigma_t(L)
 =\frac{L}{v_0^2}
 \left|\frac{\partial v}{\partial E}\right|_{E_0}\sigma_E .}
\label{eq:arrival_time_jitter_energy_spread_dp47}
\end{equation}
若第 $\Delta\ell$ 阶密度谐波的局部角频率为 $\Delta\ell\omega$，对高斯能量分布平均后得到
\begin{equation}
\boxed{
 B_{\Delta\ell}^{(E)}(L)
 =B_{\Delta\ell}^{(0)}(L)
 \exp\!\left[-\frac12\Delta\ell^2\omega^2\sigma_t^2(L)\right].}
\label{eq:bunching_energy_spread_damping_dp47}
\end{equation}
因此高阶谐波比基频更快被有限能散压低，在这个线性到达时间、高斯平均模型内，衰减指数按 $\Delta\ell^2$ 增长。这里要求 $\sigma_E$ 足够小，使速度的一阶展开有效，且随机偏移与调制振幅及环境相位独立。若同一能散同时改变散射振幅或反冲色散，就须对完整前向模型积分，而不是只乘这个相位因子。

开放系统退相干可以在同一对象上表示。若能量格的 Markov 纯退相干满足
\begin{equation}
\boxed{
 \dot\rho_{\ell\ell'}\big|_{\rm deph}
 =-\frac{\Gamma_\varphi}{2}(\ell-\ell')^2\rho_{\ell\ell'}.}
\label{eq:bunching_pure_dephasing_master_dp68}
\end{equation}
这里沿用\chapref{chap:electron-open-dynamics}的 $\Gamma_\varphi=2\int_0^\infty C_\varphi(\tau)\dd\tau$；$\mathcal D[\Lambda_e]$ 在阶差 $\Delta\ell$ 上给 $-\Delta\ell^2/2$。因此第 $\Delta\ell$ 条副对角线统一获得因子
\begin{equation}
\boxed{
 B_{\Delta\ell}^{(\mathrm{deph})}(L)
 =B_{\Delta\ell}^{(0)}(L)
 \exp\!\left[-\frac{\Delta\ell^2\Gamma_\varphi L}{2v_0}\right].}
\label{eq:bunching_dephasing_damping_dp47}
\end{equation}
能散平均与环境退相干来自不同物理机制，但都作用于同一个 $B_{\Delta\ell}$。若二者独立，实验前向模型可写成
\begin{equation}
\boxed{
 B_{\Delta\ell}^{\rm obs}(L)
 =B_{\Delta\ell}^{(0)}(L)
 \exp\!\left[-\frac12\Delta\ell^2\omega^2\sigma_t^2(L)
 -\frac{\Delta\ell^2\Gamma_\varphi L}{2v_0}\right].}
\label{eq:bunching_observed_master_dp47}
\end{equation}
这一定量式把电子源品质、漂移长度和环境相干时间直接连接到可观测密度谐波。

对给定调制强度，最优漂移长度应由最大化目标谐波而不是引用固定 Talbot 分数来确定。定义
\begin{equation}
\boxed{
 L_{\Delta\ell}^*\equiv\underset{L>0}{\operatorname{argmax}}
 \left|B_{\Delta\ell}^{\rm obs}(L)\right| .}
\label{eq:optimal_bunching_length_dp47}
\end{equation}
这里 $\operatorname{argmax}$ 可是一组距离；若上确界不取到，则不能把它记成实际最优点。在无噪声二阶色散模型中，$B_{\Delta\ell}^{(0)}(L)$ 对 $L_T$ 周期；加入\eqnrefs{eq:bunching_observed_master_dp47} 的指数包络后，不同 Talbot 周期不再等价，通常应选择最早达到局域极大的漂移距离。

若最终实验是用整形电子驱动目标跃迁，则\secref{subsec:shaped-electron-drive-dp19}的关系直接给
\begin{equation}
\boxed{
 \Omega_{3i}^{(e)}(L)
 =\frac{2G_{3i}}{T}B_{\Delta\ell_i}^{\rm obs}(L),
 \qquad
 \omega_{3i}=\Delta\ell_\ii\omega .}
\label{eq:target_rabi_with_decoherence_dp47}
\end{equation}
于是电子源能散和电子退相干不再是独立的“束流质量指标”，而是直接进入目标 Rabi 频率的乘法衰减。



\section{综合例题：同一近场模型从场积分走到探测计数}
\label{sec:feqo-nearfield-counts-comprehensive}

前面的章节分别建立了场、电子能量格、自由传播、开放演化与测量语言。
下面固定一套装置：两段经典单频近场之间插入自由漂移与相位噪声，末端用有限能量分辨率的探测器记录电子。
各步始终沿用同一个电子态，因而可以区分场的相位、电子相干性、读出脉冲与计数噪声各自的作用。
第一段和第二段的耦合记为 $\beta_1,\beta_2$；这里的 $\beta_a$ 是无量纲复耦合，不是逆温度。

\Needspace{8\baselineskip}
\begin{exbox}[第一步：从结构近场算出复耦合]
设结构产生的纵向单频近场在电子轴线上为
\[
 E_{\omega,z}(z)=E_a\exp\!\left(-\frac{z^2}{2L^2}\right)\ee^{\ii \kappa z},
 \qquad L>0.
\]
这是结构近场的轴线剖面，横向场和结构边界满足前文的 Maxwell 问题。电子以中心速度 $v_0$ 穿过近场，求其复耦合及相位失配对耦合强度的影响。

\textbf{解。} 采用前文的电荷和场幅约定，将剖面代入纵向场积分。高斯 Fourier 积分给出
\begin{align*}
 \beta&=\frac{q_e}{\hbar\omega}\int_{\mathbb R}
 E_{\omega,z}(z)\ee^{-\ii \omega z/v_0}\,\dd z\\
 &=\frac{q_e E_a\sqrt{2\pi}L}{\hbar\omega}
 \exp\!\left[-\frac{L^2}{2}\left(\kappa-\frac{\omega}{v_0}\right)^2\right].
\end{align*}
故 $\beta=g \ee^{\ii \phi}$ 中的 $g$ 随相位失配呈高斯衰减，
\[
 g=\frac{|q_e E_a|\sqrt{2\pi}L}{\hbar\omega}
 \ee^{-L^2(\kappa-\omega/v_0)^2/2},
 \qquad \phi=\arg(q_eE_a).
\]
负电子电荷带来的相位也包含在 $\phi$ 中。后文估计 $g$ 时，只有在
$L,\kappa,v_0$ 已知的条件下，才能反解场幅 $|E_a|$。
\end{exbox}

\begin{exbox}[第二步：检查能量格近似，再计算第一次散射]
确定各能量通道的动量，检查何时可用同一个复耦合计算第一次散射的振幅与能谱。

\textbf{解。} 结构可以交换动量，有限空间包络本身不指定唯一动量转移。在单频且频宽可忽略的散射模型中，先固定能量格，再由色散关系反解动量：
\[
 E_\ell=E_0+\ell\hbar\omega,\qquad
 p_\ell=\frac1c\sqrt{E_\ell^2-m^2c^4},\qquad
 k_\ell=\frac{p_{\ell+1}-p_\ell}{\hbar}.
\]
在实际占据的有限能级窗内要求 $E_\ell>mc^2$，并要求各通道的失配变化
在近场长度上不能积累显著相位，例如
\[
 \max_{|\ell|\le M}L|k_\ell-\omega/v_0|\ll1.
\]
同时须控制速度变化、场频宽及不同能级的耦合变化，使第一步的同一个 $\beta$ 可用于所有相关边。
仅有 $\hbar\omega/E_0\ll1$ 并不足以保证这些条件。

在已经满足这些条件的理想双向能量格 $\ell^2(\mathbb Z)$ 中，定义
\[
 b|\ell\rangle=|\ell-1\rangle,\qquad
 b^\dagger|\ell\rangle=|\ell+1\rangle,\qquad
 bb^\dagger=b^\dagger b=I.
\]
因此此处 $b$ 不是具有真空底态的玻色湮灭算符。近场散射为
\[
 S(\beta)=\exp(\beta b^\dagger-\beta^*b),\qquad
 S(\beta)|0\rangle=\sum_{\ell\in\mathbb Z}a_\ell|\ell\rangle,
 \qquad a_\ell=J_\ell(2g)\ee^{\ii \ell\phi}.
\]
这一步同时给出振幅与能谱：
\[
 P_\ell^{(1)}=J_\ell^2(2g),\qquad
 \sum_\ell P_\ell^{(1)}=1,
 \qquad P_0^{(1)}=1-2g^2+O(g^4),\quad
 P_{\pm1}^{(1)}=g^2+O(g^4).
\]
由 $J_{-\ell}^2=J_\ell^2$ 得平均净能级变化为零。
在相位表示中，波函数为
\[
 \psi(\theta)=\frac{1}{\sqrt{2\pi}}
 \ee^{\beta \ee^{-\ii \theta}-\beta^*\ee^{\ii \theta}},\qquad
 N=\ii\partial_\theta.
\]
Parseval 恒等式于是给出
\[
 \langle N^2\rangle
 =\int_0^{2\pi}|\partial_\theta\psi|^2\,\dd\theta
 =\frac{1}{2\pi}\int_0^{2\pi}4g^2\cos^2(\phi-\theta)\,\dd\theta
 =2g^2.
\]
电子能谱展宽了，但单次能量测量仍完全看不到 $\phi$。
\end{exbox}

\begin{exbox}[第三步：自由漂移怎样保留能谱并改变相干性]
电子穿过两近场之间的固定距离 $Z_d$。求漂移后的密度矩阵，并比较其对角元与非对角元。

\textbf{解。} 固定距离的传播相位由同一能量格的 $p(E)$ 决定。
在 $E_0$ 附近展开动量：
\[
 p_\ell=p_0+\frac{\ell\hbar\omega}{v_0}
 -\frac{(\ell\hbar\omega)^2}{2m_\parallel v_0^3}+R_\ell^{(p)},
 \qquad m_\parallel=\gamma_0^3m.
\]
在去掉中心相位与已知线性平移的参考系中，漂移时间 $T$ 对应
\[
 T=Z_d/v_0,\qquad V(Z_d)=\ee^{-\ii \varphi_rN^2},\qquad
 \varphi_r=\frac{\hbar\omega^2Z_d}{2m_\parallel v_0^3}.
\]
这里使用固定距离的传播因子 $\ee^{\ii p_\ell Z_d/\hbar}$，二次相位来自动量对能量的非线性依赖。严格等间距 $E_\ell$ 在固定时间的传播中只产生关于 $\ell$ 的线性相位。
二阶近似还要求 $\max_{|\ell|\le M}Z_d|R_\ell^{(p)}|/\hbar\ll1$。
第一次散射后的密度矩阵因而变为
\[
 \rho^{\rm drift}_{\ell m}
 =a_\ell a_m^*\ee^{-\ii \varphi_r(\ell^2-m^2)}.
\]
取 $\ell=m$ 后相位抵消，所以漂移前后的 $P_\ell^{(1)}$ 完全相同。
非对角元仍保留相位差，所以电子的纵向相干观测量会改变，即使能谱保持不变。
\end{exbox}

\begin{exbox}[第四步：同一个态通过相位噪声通道]
设这一装置的随机相位 $\chi$ 服从均值零、方差 $2D$ 的正态分布，$D>0$。求噪声平均后的密度矩阵和相位测量概率。

\textbf{解。} 对相位噪声作高斯平均，得到
\begin{align*}
 \Phi_D(\rho)&=\int_{\mathbb R}\frac{\ee^{-\chi^2/(4D)}}{\sqrt{4\pi D}}
 \ee^{-\ii \chi N}\rho \ee^{\ii \chi N}\,\dd\chi,\\
 [\Phi_D(\rho)]_{\ell m}&=\ee^{-D(\ell-m)^2}\rho_{\ell m}.
\end{align*}
这是酉共轭的概率混合，故完全正且保迹；$D=0$ 定义为恒等通道的极限。
它描述这里指定的经典相位抖动，不自动代表任意环境的微观退相干。
结合漂移，第二次散射前的同一个态为
\[
 \rho^{\rm in}_{\ell m}
 =a_\ell a_m^*\ee^{-\ii \varphi_r(\ell^2-m^2)}\ee^{-D(\ell-m)^2}.
\]
对圆周上的协变相位测量，概率密度可以写成
\[
 p(\theta)=\frac1{2\pi}\sum_{\ell,m}
 \rho^{\rm in}_{\ell m}\ee^{-\ii (\ell-m)\theta},\qquad
 \int_0^{2\pi}p(\theta)\,\dd\theta=1.
\]
保留 $a_0=1+O(g^2),a_1=\beta+O(g^3),a_{-1}=-\beta^*+O(g^3)$，得到
\[
 p(\theta)=\frac1{2\pi}
 \left[1+4g \ee^{-D}\sin\varphi_r\sin(\phi-\theta)+O(g^2)\right].
\]
第一次脉冲刚结束时的相位密度严格均匀；漂移把相位调制转成密度调制，
而噪声压低一阶谐波。这里展开式须在 $g\ll1$ 下使用，不能把其截断形式当成任意耦合下的概率密度。
\end{exbox}

\begin{exbox}[第五步：第二次近场把相干信息转成能谱]
用第二段近场读出相干信息，求最终能谱，并检验零漂移、无噪声时的相位干涉。

\textbf{解。} 读出脉冲 $S(\beta_2)$ 的矩阵元为
\[
 U_{nm}=J_{n-m}(2g_2)\ee^{\ii (n-m)\phi_2}.
\]
所以最终能谱的完整表达是
\begin{align*}
 P_n^{\rm out}
 &=\sum_{m,m'}U_{nm}U_{nm'}^*a_ma_{m'}^*
 \ee^{-\ii \varphi_r(m^2-m'^2)}\ee^{-D(m-m')^2}.
\end{align*}
双重求和保留了不同入射能级之间的相干项，其中漂移相位与噪声因子分别由前两步确定。
若 $D=0$，它就是
\[
 P_n^{\rm out}=|A_n|^2,\qquad
 A_n=\sum_mJ_{n-m}(2g_2)\ee^{\ii (n-m)\phi_2}
 \ee^{-\ii \varphi_rm^2}J_m(2g_1)\ee^{\ii m\phi_1}.
\]
特别地，零漂移且无噪声时，两个散射生成元交换，因此
\[
 S(\beta_2)S(\beta_1)=S(\beta_1+\beta_2),\qquad
 P_n^{\rm out}=J_n^2(2|\beta_1+\beta_2|).
\]
若两脉冲等强 $g_1=g_2=g$，令 $\delta=\phi_1-\phi_2$，则
\[
 |\beta_1+\beta_2|=2g|\cos(\delta/2)|,
 \qquad P_n^{\rm out}=J_n^2\!\left(4g|\cos(\delta/2)|\right).
\]
同相时脉冲增强，反相时完全抵消。这是能量测量读出相对相位的最简单极限。
存在漂移或噪声时必须回到上面的双重求和，不能继续套用这条抵消公式。
\end{exbox}

\begin{exbox}[第六步：有限分辨率与漏检写成同一组 POVM]
将有限能量分辨率和探测效率写成测量算符，求各记录窗及漏检结果的概率。

\textbf{解。} 令仪器读出的能量变化为 $\varepsilon$，高斯响应为
\[
 R(\varepsilon|n)=\frac{1}{\sqrt{2\pi}\sigma_E}
 \exp\!\left[-\frac{(\varepsilon-n\hbar\omega)^2}{2\sigma_E^2}\right],
 \qquad F_\varepsilon=\sum_nR(\varepsilon|n)|n\rangle\langle n|.
\]
因为 $\int R(\varepsilon|n)\dd\varepsilon=1$，有 $\int F_\varepsilon\dd\varepsilon=I$。
对不交叠的记录窗 $W_j$，定义
\[
 r_{jn}=\int_{W_j}R(\varepsilon|n)\,\dd\varepsilon,\qquad
 F_j=\eta_j\sum_nr_{jn}|n\rangle\langle n|,
 \quad 0\le\eta_j\le1.
\]
于是记录概率为
\[
 p_j=\operatorname{Tr}(F_j\rho^{\rm out})
 =\eta_j\sum_nr_{jn}P_n^{\rm out}.
\]
有限窗口或效率低于一时，必须补上
\[
 F_{\rm miss}=I-\sum_jF_j\ge0,\qquad
 p_{\rm miss}=1-\sum_jp_j.
\]
这一步把能级混叠与漏检分开了。$\sigma_E\ll\hbar\omega$ 时可以近似分辨边带；
当窗口混叠明显时，真实边带占据不能直接等同于该窗口计数比例。
\end{exbox}

\begin{exbox}[第七步：计数似然、相位可辨识性与一个数值设计]
由记录概率构造计数似然，判断单个参考相位设置能否确定相对相位。

\textbf{解。} 给定固定入射电子数 $N_{\rm inc}$，把漏检也作为结果，全部计数服从多项分布。
这里要求各电子独立且同态制备，相位噪声按电子独立采样。
若一批电子共享同一个慢变化相位，单电子通道仍可相同，批计数却是条件分布的混合，不能直接沿用独立分流。
若入射数本身为均值 $\Lambda$ 的泊松变量，各窗口分流后得到独立泊松计数。再加上相互独立、已标定均值 $B_j$ 的泊松背景，有
\[
 C_j\sim\operatorname{Poisson}(\lambda_j),\qquad
 \lambda_j(\delta)=\Lambda p_j(\delta)+B_j.
\]
似然、得分和 Fisher 信息依次为
\begin{align*}
 \log\mathcal L(\delta)&=\sum_j[C_j\log\lambda_j-\lambda_j-\log(C_j!)],\\
 \partial_\delta\log\mathcal L&=\sum_j
 \left(\frac{C_j}{\lambda_j}-1\right)\partial_\delta\lambda_j,\\
 \mathcal I(\delta)&=\sum_{j:\lambda_j>0}
 \frac{(\partial_\delta\lambda_j)^2}{\lambda_j}.
\end{align*}
得分与通常的 Fisher 公式在所用窗口均有 $\lambda_j>0$ 的正则参数邻域成立。
理想反相、无背景的暗窗口满足 $\lambda_j=0$，此时应直接使用原始似然及相应极限；
不能将 $0/0$ 代入，也不能因删去该项就宣布暗端没有参数信息。
只有第一次脉冲并直接测能量时，$p_j$ 与 $\phi_1$ 无关，故这项相位信息严格为零。
读出脉冲引入后，零漂移、无噪声、等强弱脉冲的两条一阶边带合概率为
\[
 q(\delta)=P_1^{\rm out}+P_{-1}^{\rm out}
 =4g^2(1+\cos\delta)+O(g^4).
\]
一次读出仍只看到 $\cos\delta$，不能区分 $\delta$ 与 $-\delta$。
将已知参考相位增加 $\pi/2$ 后重复实验，两个设置分别测得余弦与正弦分量，
才可以在模 $2\pi$ 意义下消除这项歧义。

取 $g=0.10$、每个设置 $N_{\rm inc}=10^4$、理想分辨率、相同效率 $\eta=0.8$，忽略背景。
按弱耦合近似，$\delta=0,\pi/2,\pi$ 的两边带期望计数依次为
\[
 \eta N_{\rm inc}q(\delta)\simeq640,\quad320,\quad0.
\]
这是设计期望，并非虚构的实测数据。精确预测应使用 Bessel 概率；
其中反相设置在当前理想极限下严格为零，另两个数含弱耦合截断误差。
\end{exbox}

\begin{exbox}[第八步：反演场幅与控制计算窗]
由第一次散射的计数估计场幅，并用边带尾部估计有限计算窗的截断误差。

\textbf{解。} 先关闭读出脉冲，对第一次散射的两条一阶边带做理想分辨率标定。
在弱耦合、已知效率 $\eta$、无背景条件下，边带总计数 $C$ 近似满足
\[
 C\sim\operatorname{Binomial}(N_{\rm inc},2\eta g^2),\qquad
 \widehat g=\sqrt{\frac{C}{2\eta N_{\rm inc}}}.
\]
它是弱耦合模型的反演估计，并非精确 Bessel 模型下的无偏估计。
当期望计数足够大，Delta 方法给出
\[
 \operatorname{Var}(\widehat g)
 \simeq\frac{1-2\eta g^2}{8\eta N_{\rm inc}},\qquad
 |\widehat E_a|=\frac{\hbar\omega\widehat g}{|q_e|\sqrt{2\pi}L}
 \ee^{L^2(\kappa-\omega/v_0)^2/2}.
\]
因此大的相位失配既压低信号，也放大场幅反演的不确定性。
若有背景、窗口混叠或未知效率，就必须使用第六、七步的完整概率模型联合拟合；
不能只减去平均背景后仍沿用无背景二项方差。

对第一次散射，第二步算出的二阶矩还给出可直接用于数值截断的界
\[
 \epsilon_M=\sum_{|\ell|>M}|a_\ell|^2
 \le\frac{2g^2}{(M+1)^2}.
\]
归一化截断纯态与原态的迹距离为 $\sqrt{\epsilon_M}$，
故在后续使用相同的完全正保迹演化和相同测量时，结果概率的总变差不超过这一界。
这不包括把第二次脉冲自身替换为有限矩阵所增加的边界误差；
也不自动控制能量这类无界观测量的期望，后者需要另行估计尾部矩。

\end{exbox}

场剖面的 $E_a,L,\kappa,v_0$ 先决定 $\beta_1$，随后得到的电子振幅经历漂移、噪声和参考脉冲，最终转成记录概率与计数似然。因而场幅反演依赖整条演化和测量链。色散截断、相位噪声模型、仪器响应和计数采样各有自己的使用条件；改变其中一项，就应从对应的那一步重新计算预测与不确定性。

